% ============================================================
% ANEXO — FICHAMENTOS SELECIONADOS
% ============================================================

\chapter{Fichamentos Selecionados}
\label{anexo:fichamentos}

Este anexo reúne os fichamentos dos principais textos que fundamentam a presente pesquisa. Cada fichamento segue o protocolo estabelecido em \texttt{fichamentos/README.md}, contendo a referência completa, os conceitos-chave extraídos, as passagens mais relevantes e a articulação com o problema de pesquisa.

\section{Magistério Pontifício}

\begin{fichamento}
    \textbf{Ref.:} \cite{LeaoXIV2026MH}
    
    \textbf{Conceitos-chave:} dignidade humana, paradigma tecnocrático, irredutibilidade da pessoa, desenvolvimento integral, ecologia integral, destino universal dos bens, subsidiariedade, solidariedade.
    
    \textbf{Passagens relevantes:}
    \begin{itemize}
        \item §3: ``A pessoa humana [...] não pode ser reduzida à soma de seus dados, a um perfil calculável ou a um conjunto de preferências observáveis.''
        \item §89: ``A tecnologia, expressão da criatividade humana, não é em si mesma boa nem má; tudo depende da orientação ética que lhe é dada.''
        \item §101: ``A irredutibilidade da pessoa é o primeiro princípio de toda regulação que pretenda proteger a dignidade humana na era digital.''
        \item §186: ``A pessoa não pode ser transformada em objeto de classificação algorítmica.''
        \item §232: ``Toda a técnica deve estar a serviço do desenvolvimento integral da pessoa humana.''
    \end{itemize}
    
    \textbf{Articulação:} Documento central da pesquisa. Oferece o referencial teórico e principiológico para a análise da proteção da personalidade na era da IA.
\end{fichamento}

\begin{fichamento}
    \textbf{Ref.:} \cite{Francisco2015LS}
    
    \textbf{Conceitos-chave:} ecologia integral, paradigma tecnocrático, cuidado da casa comum, destino universal dos bens, diálogo intergeracional.
    
    \textbf{Passagens relevantes:}
    \begin{itemize}
        \item §106–114: Crítica ao paradigma tecnocrático como visão de mundo que reduz a realidade a objeto de manipulação.
        \item §115: Identificação da raiz humana da crise ecológica.
        \item §203–204: Exigência de uma ecologia integral que articule a dimensão ambiental, social e humana.
    \end{itemize}
    
    \textbf{Articulação:} A crítica ao paradigma tecnocrático presente na \textit{Laudato Si'} é retomada e atualizada pela \MH{} no contexto específico da IA. Constitui a base teórica para a análise crítica do capítulo quinto.
\end{fichamento}

\begin{fichamento}
    \textbf{Ref.:} \cite{JoaoXXIII1963PT}
    
    \textbf{Conceitos-chave:} bem comum, dignidade humana, autoridade política, participação, subsidiariedade.
    
    \textbf{Passagens relevantes:}
    \begin{itemize}
        \item §8–9: Definição do bem comum como conjunto de condições sociais que permitem o desenvolvimento integral da pessoa.
        \item §51–52: A autoridade política como serviço ao bem comum.
        \item §64–65: O princípio da subsidiariedade como limite à atuação do Estado.
    \end{itemize}
    
    \textbf{Articulação:} A doutrina social da Igreja sobre o bem comum e a subsidiariedade, desenvolvida na \textit{Pacem in Terris}, fundamenta os princípios éticos que a \MH{} aplica à regulação da IA.
\end{fichamento}

\begin{fichamento}
    \textbf{Ref.:} \cite{JoaoPauloII1991CA}
    
    \textbf{Conceitos-chave:} livre iniciativa, destino universal dos bens, subsidiariedade, solidariedade, Estado regulador.
    
    \textbf{Passagens relevantes:}
    \begin{itemize}
        \item §15: Crítica ao capitalismo desregulado.
        \item §48: O Estado como promotor do desenvolvimento, não mero espectador.
    \end{itemize}
    
    \textbf{Articulação:} Os princípios da subsidiariedade e da solidariedade, articulados na \textit{Centesimus Annus}, são retomados pela \MH{} como critérios para a governança algorítmica.
\end{fichamento}

\begin{fichamento}
    \textbf{Ref.:} \cite{Francisco2020FT}
    
    \textbf{Conceitos-chave:} fraternidade universal, amizade social, diálogo intergeracional, bem comum, cultura do encontro, solidariedade.
    
    \textbf{Passagens relevantes:}
    \begin{itemize}
        \item §20–21: ``O sonho de fraternidade'' que fundamenta a convivência humana para além de qualquer fronteira.
        \item §34–36: O amor universal como princípio ético que reconhece cada pessoa como irmã, independentemente de sua origem ou condição.
        \item §55: A regra de ouro — não fazer ao outro o que não queres que te façam — como princípio fundamental do diálogo intercultural.
        \item §101–104: A amizade social como condição para a construção de um mundo mais justo e solidário.
        \item §205: Crítica ao individualismo contemporâneo e apelo à solidariedade como virtude política fundamental.
    \end{itemize}
    
    \textbf{Articulação:} A fraternidade universal proposta pela \textit{Fratelli Tutti} oferece um princípio ético fundamental para orientar a governança da IA, assegurando que o desenvolvimento tecnológico esteja a serviço da solidariedade e do bem comum, conforme retomado pela \MH{} no contexto digital.
\end{fichamento}

\begin{fichamento}
    \textbf{Ref.:} \cite{ConcilioVaticanoII1965GS}
    
    \textbf{Conceitos-chave:} dignidade da pessoa humana, progresso técnico, bem comum, autonomia das realidades terrenas, relação Igreja-mundo, desenvolvimento integral.
    
    \textbf{Passagens relevantes:}
    \begin{itemize}
        \item §26: O bem comum como ``conjunto das condições da vida social que permitem aos grupos e a cada um dos seus membros atingir de modo mais completo e facilmente a própria perfeição.''
        \item §34: O progresso técnico como expressão da criatividade humana, que deve estar ordenado ao serviço da pessoa.
        \item §36: A autonomia legítima das realidades terrenas — a ciência e a técnica têm suas leis próprias, mas devem submeter-se aos valores morais fundamentais.
        \item §63: As transformações sociais decorrentes do progresso técnico e seus desafios éticos.
        \item §76: A autonomia da ordem temporal e sua relação com a missão da Igreja na promoção da dignidade humana.
    \end{itemize}
    
    \textbf{Articulação:} A Constituição Pastoral \textit{Gaudium et Spes} estabelece o princípio da autonomia legítima das realidades terrenas (§36) que a \MH{} retoma ao reconhecer o valor intrínseco da técnica, mas também a necessidade de subordiná-la aos imperativos éticos fundamentais, especialmente a dignidade da pessoa humana diante do reducionismo algorítmico.
\end{fichamento}

\section{Livros e Capítulos de Livros}

\begin{fichamento}
    \textbf{Ref.:} \cite{Bioni2021PPD}
    
    \textbf{Conceitos-chave:} proteção de dados pessoais, autodeterminação informativa, privacidade, LGPD, direito à explicação.
    
    \textbf{Conceitos principais:}
    \begin{itemize}
        \item A autodeterminação informativa como fundamento da proteção de dados.
        \item O direito à explicação como garantia do titular diante de decisões automatizadas.
        \item A LGPD como marco regulatório brasileiro da proteção de dados.
    \end{itemize}
    
    \textbf{Articulação:} Fundamenta a análise crítica da LGPD no capítulo quarto e a articulação com os princípios da \MH{} no capítulo quinto.
\end{fichamento}

\begin{fichamento}
    \textbf{Ref.:} \cite{Doneda2006PDD}
    
    \textbf{Conceitos-chave:} proteção de dados, privacidade, direito comparado, evolução legislativa.
    
    \textbf{Conceitos principais:}
    \begin{itemize}
        \item A proteção de dados como direito autônomo.
        \item A evolução legislativa europeia e brasileira.
        \item A necessidade de harmonização entre proteção de dados e direitos da personalidade.
    \end{itemize}
    
    \textbf{Articulação:} O marco teórico da proteção de dados desenvolvido por Doneda fornece as bases conceituais para a análise do regime jurídico brasileiro no capítulo quarto.
\end{fichamento}

\begin{fichamento}
    \textbf{Ref.:} \cite{Rodota2008PDA}
    
    \textbf{Conceitos-chave:} privacidade, proteção de dados, dignidade, identidade pessoal, sociedade informacional.
    
    \textbf{Conceitos principais:}
    \begin{itemize}
        \item A privacidade como condição para o exercício da liberdade.
        \item A proteção de dados como instrumento de tutela da identidade pessoal.
        \item O direito de controlar a própria informação como dimensão da dignidade.
    \end{itemize}
    
    \textbf{Articulação:} A conexão entre proteção de dados e dignidade humana, sustentada por Rodotà, dialoga diretamente com a fundamentação antropológica da \MH{}, especialmente a tese da irredutibilidade da pessoa aos seus dados.
\end{fichamento}

\begin{fichamento}
    \textbf{Ref.:} \cite{Castells2010EI}
    
    \textbf{Conceitos-chave:} sociedade em rede, era da informação, poder, identidade, globalização.
    
    \textbf{Conceitos principais:}
    \begin{itemize}
        \item A sociedade em rede como nova estrutura social da era da informação.
        \item O poder dos fluxos de informação como nova forma de dominação.
        \item A crise da identidade pessoal na sociedade em rede.
    \end{itemize}
    
    \textbf{Articulação:} O quadro sociológico de Castells fornece o contexto para compreender os desafios à personalidade humana na era digital, examinados no capítulo terceiro.
\end{fichamento}

\begin{fichamento}
    \textbf{Ref.:} \cite{Zuboff2019SC}
    
    \textbf{Conceitos-chave:} capitalismo de vigilância, excedente comportamental, sociedade instrumental, direitos humanos.
    
    \textbf{Conceitos principais:}
    \begin{itemize}
        \item O capitalismo de vigilância como nova ordem econômica que trata a experiência humana como matéria-prima gratuita.
        \item A instrumentalização da pessoa como ameaça aos direitos humanos.
        \item A necessidade de um novo contrato social para a era digital.
    \end{itemize}
    
    \textbf{Articulação:} A teoria do capitalismo de vigilância fundamenta a análise dos riscos da IA à personalidade no capítulo terceiro e dialoga com a crítica da \MH{} ao paradigma tecnocrático no capítulo quinto.
\end{fichamento}

\begin{fichamento}
    \textbf{Ref.:} \cite{AquinoST}
    
    \textbf{Conceitos-chave:} pessoa, dignidade, lei natural, bem comum, substância individual de natureza racional, sindérese, virtude.
    
    \textbf{Conceitos principais:}
    \begin{itemize}
        \item A definição de pessoa como ``substância individual de natureza racional'' (\textit{individua substantia rationalis naturae}), que estabelece a dignidade ontológica do ser humano (I, q. 29, a. 1).
        \item A lei natural como participação da lei eterna na criatura racional, fundamento objetivo da moralidade (I-II, q. 91, a. 2).
        \item O bem comum como fim da sociedade política e critério de justiça das leis positivas (I-II, q. 90, a. 2).
        \item A sindérese como hábito natural dos primeiros princípios práticos, base da consciência moral (I, q. 79, a. 12).
        \item A virtude como aperfeiçoamento das faculdades humanas para o agir segundo a razão reta (I-II, q. 55, a. 1).
    \end{itemize}
    
    \textbf{Articulação:} A antropologia filosófica de Tomás de Aquino, em particular a definição de pessoa como substância individual de natureza racional, fundamenta a tese central da \MH{} sobre a irredutibilidade da pessoa face aos sistemas de IA. O conceito de lei natural informa a crítica ao paradigma tecnocrático, que subordina o bem comum à lógica algorítmica, enquanto a doutrina do bem comum fornece o critério teleológico para a regulação ética da inteligência artificial.
\end{fichamento}

\begin{fichamento}
    \textbf{Ref.:} \cite{Mounier1949P}
    
    \textbf{Conceitos-chave:} pessoa, engajamento, comunidade, liberdade, personalismo comunitário, transcendência, vocação.
    
    \textbf{Conceitos principais:}
    \begin{itemize}
        \item A pessoa como ``ser espiritual, naturalmente independente, transcendente às realidades materiais e sociais, mas que se realiza na comunhão com as outras pessoas''.
        \item A distinção entre individualidade (fechamento egoísta) e personalidade (abertura ao outro e à comunidade).
        \item O engajamento pessoal como condição da liberdade autêntica — a pessoa se realiza na ação responsável.
        \item A comunidade como meio e fim do desenvolvimento pessoal, contraposta ao individualismo liberal e ao coletivismo totalitário.
        \item A transcendência como dimensão constitutiva da pessoa que a abre ao absoluto e ao infinito.
    \end{itemize}
    
    \textbf{Articulação:} O personalismo comunitário de Mounier fornece o quadro filosófico que informa a concepção de pessoa na \MH{}. A crítica mounieriana à redução da pessoa a funções sociais ou econômicas antecipa a crítica da encíclica à redução algorítmica da pessoa a perfis de dados. A comunidade como realização da pessoa orienta a visão de uma governança algorítmica participativa e solidária.
\end{fichamento}

\begin{fichamento}
    \textbf{Ref.:} \cite{Ricoeur1990SM}
    
    \textbf{Conceitos-chave:} identidade narrativa, ipseidade, alteridade, estima de si, solicitude, sabedoria prática.
    
    \textbf{Conceitos principais:}
    \begin{itemize}
        \item A distinção entre identidade-idem (\textit{idem} — mesmidade, caráter fixo) e identidade-ipse (\textit{ipse} — ipseidade, si-mesmo constituído na promessa e na narrativa).
        \item A ipseidade como identidade dinâmica que se constitui na relação com o outro e na fidelidade à palavra dada.
        \item A alteridade constitutiva do si-mesmo — ``si-mesmo como um outro'' —, que insere a dimensão ética no núcleo da identidade pessoal.
        \item A estima de si e a solicitude como mediações éticas da identidade narrativa.
        \item A sabedoria prática como julgamento em situação concreta, para além da aplicação mecânica de regras universais.
    \end{itemize}
    
    \textbf{Articulação:} A noção de identidade narrativa de Ricoeur fundamenta a crítica à redução da pessoa a perfis algorítmicos na era da IA. Se a identidade pessoal se constitui narrativamente, através de um processo contínuo de interpretação e reinterpretação, então os sistemas de IA que classificam, predizem e definem a pessoa com base em dados passados cometem um erro antropológico fundamental: confundem a mesmidade (identidade-idem) com a totalidade da pessoa, suprimindo a dimensão de abertura e novidade própria da ipseidade.
\end{fichamento}

\begin{fichamento}
    \textbf{Ref.:} \cite{Jonas1979PR}
    
    \textbf{Conceitos-chave:} responsabilidade, ética do futuro, imperativo da responsabilidade, heurística do medo, vulnerabilidade, precaução.
    
    \textbf{Conceitos principais:}
    \begin{itemize}
        \item O imperativo da responsabilidade: ``Age de tal modo que os efeitos da tua ação sejam compatíveis com a permanência de uma vida humana autêntica sobre a Terra.''
        \item A heurística do medo (\textit{Heuristik der Furcht}) como método para antecipar cenários catastróficos antes que se concretizem.
        \item A vulnerabilidade da natureza e das gerações futuras como fundamento de uma nova ética da responsabilidade.
        \item A crítica ao utopismo tecnológico que ignora os riscos de longo prazo do desenvolvimento técnico.
        \item A assimetria entre o poder da técnica e a capacidade humana de prever suas consequências.
    \end{itemize}
    
    \textbf{Articulação:} A ética da responsabilidade de Jonas fundamenta o princípio de precaução aplicado à IA, retomado pela \MH{} na crítica ao paradigma tecnocrático. A heurística do medo oferece um método para identificar riscos existenciais da IA antes que se materializem, enquanto o imperativo da responsabilidade estabelece o dever ético de assegurar que o desenvolvimento algorítmico seja compatível com a dignidade e a autonomia da pessoa humana, orientando a regulação prudencial proposta no capítulo quinto.
\end{fichamento}

\begin{fichamento}
    \textbf{Ref.:} \cite{Floridi2023EI}
    
    \textbf{Conceitos-chave:} ética da IA, princípios éticos, governança algorítmica, dignidade informacional, sociedade da informação, infraética.
    
    \textbf{Conceitos principais:}
    \begin{itemize}
        \item A ética da IA como campo disciplinar autônomo que articula princípios normativos para o desenvolvimento e uso de sistemas inteligentes.
        \item A distinção entre ética aplicada (\textit{applied ethics}) e infraética (\textit{infraethics}), que fornece o enquadramento normativo de fundo para a governança algorítmica.
        \item A dignidade informacional como princípio ético fundamental que reconhece o valor intrínseco de cada pessoa na sociedade da informação.
        \item Os quatro princípios éticos para a IA: beneficência, não maleficência, autonomia e justiça, acrescidos da explicabilidade como princípio transversal.
        \item A necessidade de traduzir princípios éticos em práticas regulatórias efetivas, superando a mera declaração de intenções.
    \end{itemize}
    
    \textbf{Articulação:} Oferece o marco ético contemporâneo da IA que dialoga com os princípios propostos pela \MH{}. A dignidade informacional de Floridi converge com a irredutibilidade da pessoa sustentada pela encíclica, enquanto os princípios éticos propostos pelo autor fornecem um quadro comparativo para os princípios da \MH{} examinados no capítulo segundo, permitindo uma análise crítica dos limites e potencialidades de cada abordagem.
\end{fichamento}

\begin{fichamento}
    \textbf{Ref.:} \cite{MendesDoneda2021PD}
    
    \textbf{Conceitos-chave:} proteção de dados, LGPD, direito comparado, dignidade, autodeterminação informativa, privacidade, ANPD.
    
    \textbf{Conceitos principais:}
    \begin{itemize}
        \item A proteção de dados pessoais como direito fundamental autônomo, para além da mera tutela da privacidade.
        \item A análise crítica da LGPD: avanços e insuficiências do marco regulatório brasileiro diante dos desafios tecnológicos.
        \item O direito à explicação como garantia específica diante de decisões automatizadas e seus limites na legislação brasileira.
        \item A Autoridade Nacional de Proteção de Dados (ANPD) como instituição central para a efetividade do sistema de proteção.
        \item A perspectiva comparada com o direito europeu e norte-americano como método para identificar lacunas e boas práticas.
    \end{itemize}
    
    \textbf{Articulação:} A perspectiva crítica sobre a proteção de dados informa a análise das insuficiências do regime brasileiro diante dos desafios da IA, desenvolvida no capítulo quarto. A ênfase na proteção de dados como direito fundamental autônomo dialoga com a fundamentação antropológica da \MH{}, que reconhece na dignidade da pessoa o fundamento último de toda proteção jurídica na era digital.
\end{fichamento}

\begin{fichamento}
    \textbf{Ref.:} \cite{Comparato2016DH}
    
    \textbf{Conceitos-chave:} direitos humanos, dignidade, evolução histórica, universalidade, indivisibilidade, fundamentação filosófica.
    
    \textbf{Conceitos principais:}
    \begin{itemize}
        \item A afirmação histórica dos direitos humanos como processo de positivação da dignidade da pessoa humana ao longo dos séculos.
        \item A universalidade dos direitos humanos como expressão da igual dignidade de todos os seres humanos, independentemente de sua condição.
        \item A indivisibilidade dos direitos humanos — direitos civis, políticos, econômicos, sociais e culturais formam um conjunto unitário.
        \item A fundamentação filosófica dos direitos humanos na tradição do direito natural e no personalismo ético.
        \item O reconhecimento dos novos direitos diante das transformações técnicas e sociais como característica dinâmica do sistema de direitos humanos.
    \end{itemize}
    
    \textbf{Articulação:} A fundamentação histórica dos direitos humanos desenvolvida por Comparato fornece o pano de fundo para compreender a proteção da personalidade como direito fundamental em evolução. A tese do autor sobre a abertura do sistema de direitos humanos a novas exigências éticas fundamenta a proposta de extensão da tutela da personalidade diante dos desafios específicos da IA, em consonância com a visão da \MH{} sobre o desenvolvimento integral.
\end{fichamento}

\begin{fichamento}
    \textbf{Ref.:} \cite{Sarlet2021DF}
    
    \textbf{Conceitos-chave:} direitos fundamentais, dimensão objetiva, eficácia horizontal, dignidade, proporcionalidade, restrição a direitos.
    
    \textbf{Conceitos principais:}
    \begin{itemize}
        \item A dupla dimensão dos direitos fundamentais: subjetiva (direitos subjetivos) e objetiva (valores objetivos que informam todo o ordenamento).
        \item A eficácia horizontal dos direitos fundamentais como vinculação dos particulares, especialmente relevante para as relações entre plataformas digitais e usuários.
        \item A dignidade da pessoa humana como fundamento e limite de todos os direitos fundamentais, com eficácia irradiante sobre todo o ordenamento jurídico.
        \item O princípio da proporcionalidade como método de solução de conflitos entre direitos fundamentais na era digital.
        \item A proteção contra automações — o direito a não ser submetido a decisões exclusivamente automatizadas — como nova dimensão dos direitos fundamentais.
    \end{itemize}
    
    \textbf{Articulação:} A dogmática dos direitos fundamentais informa a análise da eficácia dos direitos da personalidade diante dos desafios da IA, especialmente no que tange à eficácia horizontal (vinculação de plataformas privadas) e à dupla dimensão objetiva e subjetiva. Esta análise fundamenta a tese, desenvolvida nos capítulos quarto e quinto, de que a proteção da personalidade na era da IA exige tanto mecanismos subjetivos de tutela quanto a adoção de políticas públicas orientadas pelos valores objetivos da ordem constitucional.
\end{fichamento}

\begin{fichamento}
    \textbf{Ref.:} \cite{Frazao2020RD}
    
    \textbf{Conceitos-chave:} digitalização, riscos, regulação, responsabilidade civil, plataformas, assimetria de poder, proteção do consumidor.
    
    \textbf{Conceitos principais:}
    \begin{itemize}
        \item A digitalização como fenômeno jurídico multifacetado que impõe novos riscos ao direito privado e à proteção da pessoa.
        \item A assimetria de poder entre plataformas digitais e usuários como desafio central para a regulação do ambiente digital.
        \item A responsabilidade civil dos provedores de serviços algorítmicos diante de danos à personalidade.
        \item A insuficiência dos modelos tradicionais de responsabilidade baseados na culpa diante dos riscos sistêmicos da IA.
        \item A necessidade de uma abordagem regulatória baseada em risco, que considere o potencial danoso dos sistemas algorítmicos.
    \end{itemize}
    
    \textbf{Articulação:} A análise dos riscos da digitalização desenvolvida por Frazão fundamenta a identificação dos desafios à proteção da personalidade na era da IA, examinados no capítulo terceiro. Sua proposta de regulação baseada em risco dialoga diretamente com a abordagem da \MH{}, que também propõe uma avaliação ética prévia dos sistemas tecnológicos como condição para sua adoção legítima.
\end{fichamento}

\begin{fichamento}
    \textbf{Ref.:} \cite{Pasquale2015BM}
    
    \textbf{Conceitos-chave:} opacidade algorítmica, caixa-preta, assimetria de poder, regulação, accountability, sociedade da informação.
    
    \textbf{Conceitos principais:}
    \begin{itemize}
        \item A opacidade dos algoritmos como problema de governança democrática — decisões que afetam a vida das pessoas são tomadas por sistemas incompreensíveis.
        \item A assimetria de poder entre as corporações tecnológicas e os cidadãos na era dos dados massivos.
        \item A necessidade de transparência algorítmica como condição para a accountability e o controle social dos sistemas automatizados.
        \item O ciclo vicioso entre coleta de dados, profiling e discriminação como ameaça sistêmica à autonomia pessoal.
        \item A proposta de uma ``sociedade da caixa-preta aberta'' (\textit{open black box society}) como ideal regulatório.
    \end{itemize}
    
    \textbf{Articulação:} A crítica à opacidade dos sistemas algorítmicos desenvolvida por Pasquale dialoga com a exigência de transparência da \MH{} para a governança da IA. O conceito de ``sociedade da caixa-preta'' fundamenta a análise crítica do capítulo terceiro sobre os riscos à personalidade, enquanto a proposta de abertura dos sistemas algorítmicos informa as recomendações regulatórias do capítulo quinto, em consonância com o princípio de subsidiariedade defendido pela encíclica.
\end{fichamento}

\begin{fichamento}
    \textbf{Ref.:} \cite{Amaral2003DBC}
    
    \textbf{Conceitos-chave:} pessoa, capacidade, personalidade, direitos da personalidade, bem comum, sujeito de direito.
    
    \textbf{Conceitos principais:}
    \begin{itemize}
        \item A pessoa natural como sujeito de direito — capacidade jurídica como atributo de todo ser humano.
        \item Os direitos da personalidade como direitos inatos, absolutos, intransmissíveis, irrenunciáveis e imprescritíveis.
        \item A distinção entre capacidade de direito (aptidão para ser titular de direitos) e capacidade de fato (aptidão para exercê-los pessoalmente).
        \item O nome, a imagem, a honra, a privacidade e a integridade moral como projeções da personalidade tuteladas pelo direito civil.
        \item A teoria geral dos direitos da personalidade como construção dogmática aberta a novas manifestações da personalidade.
    \end{itemize}
    
    \textbf{Articulação:} A dogmática civilística dos direitos da personalidade, sistematizada por Amaral, fornece o quadro jurídico fundamental para a pesquisa. A abertura do sistema de direitos da personalidade a novas manifestações — como o direito à autodeterminação informativa e à explicação algorítmica — fundamenta a tese, defendida no capítulo quarto, de que a proteção contra os riscos da IA constitui uma nova dimensão dos direitos da personalidade, compatível com a visão integral da pessoa humana proposta pela \MH{}.
\end{fichamento}

\begin{fichamento}
    \textbf{Ref.:} \cite{Noble2018AFS}
    
    \textbf{Conceitos-chave:} discriminação algorítmica, racismo algorítmico, vieses, busca online, desigualdade, justiça social.
    
    \textbf{Conceitos principais:}
    \begin{itemize}
        \item A discriminação algorítmica como reprodução e amplificação das desigualdades sociais existentes através de sistemas automatizados.
        \item O racismo algorítmico — como os sistemas de busca e recomendação perpetuam estereótipos raciais e de gênero.
        \item A opacidade dos critérios de ranqueamento e recomendação como obstáculo à identificação e correção de vieses discriminatórios.
        \item A responsabilidade social das empresas de tecnologia na construção de sistemas equitativos.
        \item A necessidade de intervenção regulatória para assegurar que os sistemas algorítmicos respeitem princípios de justiça e não discriminação.
    \end{itemize}
    
    \textbf{Articulação:} A análise das discriminações algorítmicas desenvolvida por Noble fundamenta a necessidade de regulação da IA que a \MH{} defende. O estudo dos vieses algorítmicos na busca online ilustra uma das ameaças concretas à personalidade examinadas no capítulo terceiro — a objetificação algorítmica que reduz a pessoa a estereótipos — e reforça o argumento da encíclica sobre a necessidade de submeter a técnica a critérios éticos que respeitem a igual dignidade de todos os seres humanos.
\end{fichamento}

\begin{fichamento}
    \textbf{Ref.:} \cite{Pariser2022THP}
    
    \textbf{Conceitos-chave:} bolha de filtros, personalização, manipulação, autonomia, economia da atenção, fragmentação social.
    
    \textbf{Conceitos principais:}
    \begin{itemize}
        \item A bolha de filtros (\textit{filter bubble}) como isolamento informacional causado pela personalização algorítmica do conteúdo online.
        \item A personalização algorítmica como ameaça à autonomia pessoal — o usuário não escolhe o que vê, mas é conduzido por critérios opacos.
        \item A fragmentação do espaço público em bolhas incomunicantes como risco à democracia e ao diálogo social.
        \item A economia da atenção como motor dos sistemas de personalização — o conteúdo é selecionado para maximizar o engajamento, não para informar ou formar.
        \item A perda de controle sobre o próprio ambiente informacional como dimensão da crise da autonomia na era digital.
    \end{itemize}
    
    \textbf{Articulação:} O fenômeno das bolhas informacionais identificado por Pariser ilustra a ameaça à autonomia pessoal que a \MH{} denuncia no paradigma digital. A fragmentação do espaço público pela personalização algorítmica contradiz a vocação comunitária e dialógica da pessoa humana afirmada pela encíclica, enquanto a crítica à economia da atenção fundamenta a necessidade de uma regulação ética dos sistemas algorítmicos de recomendação.
\end{fichamento}

\begin{fichamento}
    \textbf{Ref.:} \cite{Kurzweil2005THS}
    
    \textbf{Conceitos-chave:} singularidade tecnológica, transumanismo, superinteligência, imortalidade, convergência NBIC, pós-humanismo.
    
    \textbf{Conceitos principais:}
    \begin{itemize}
        \item A lei dos retornos acelerados — o progresso tecnológico segue uma trajetória exponencial que levará à singularidade por volta de 2045.
        \item A singularidade tecnológica como momento em que a inteligência artificial superará a inteligência humana em todas as dimensões.
        \item O transumanismo como projeto de superação dos limites biológicos humanos através da fusão com máquinas inteligentes.
        \item A imortalidade técnica como objetivo último do projeto transumanista — o upload da mente humana para suportes digitais.
        \item A convergência NBIC (nanotecnologia, biotecnologia, tecnologia da informação e ciência cognitiva) como base material da transformação pós-humana.
    \end{itemize}
    
    \textbf{Articulação:} O transumanismo de Kurzweil representa a visão que a \MH{} critica como reducionismo tecnológico da pessoa humana. A tese da singularidade tecnológica e da superação dos limites biológicos contradiz a antropologia integral da encíclica, que afirma a irredutibilidade da pessoa à sua dimensão meramente biológica ou informacional. A análise do transumanismo no capítulo terceiro permite demonstrar que o projeto de fusão homem-máquina representa a forma mais radical de objetificação algorítmica, contra a qual a \MH{} oferece uma alternativa fundada na dignidade transcendente da pessoa.
\end{fichamento}

\begin{fichamento}
    \textbf{Ref.:} \cite{Sarmento2016DPF}
    
    \textbf{Conceitos-chave:} dignidade, mínimo existencial, autonomia, reconhecimento, vulnerabilidade, igualdade.
    
    \textbf{Conceitos principais:}
    \begin{itemize}
        \item A dignidade da pessoa humana como princípio constitucional fundamental que irradia efeitos sobre todo o ordenamento jurídico.
        \item A dupla dimensão da dignidade: negativa (não instrumentalização) e positiva (promoção das condições de vida digna).
        \item O mínimo existencial como núcleo intangível da dignidade que inclui tanto condições materiais quanto elementos imateriais (reconhecimento, autonomia, integridade moral).
        \item A autonomia pessoal como expressão da dignidade — capacidade de autodeterminação que deve ser protegida inclusive contra ingerências tecnológicas.
        \item O reconhecimento intersubjetivo como dimensão essencial da dignidade — a pessoa se realiza na relação com o outro e precisa ser reconhecida em seu valor intrínseco.
    \end{itemize}
    
    \textbf{Articulação:} A fundamentação da dignidade como direito fundamental desenvolvida por Sarmento informa a crítica da \MH{} à instrumentalização algorítmica da pessoa. A dupla dimensão da dignidade (negativa e positiva) fornece o quadro para compreender tanto as violações diretas da personalidade pela IA (vigilância, discriminação) quanto as omissões estatais na promoção de condições de autodeterminação digital. O reconhecimento como elemento da dignidade fundamenta a exigência de que os sistemas algorítmicos tratem cada pessoa como fim, jamais como meio.
\end{fichamento}

\begin{fichamento}
    \textbf{Ref.:} \cite{Doneda2022DPD}
    
    \textbf{Conceitos-chave:} direitos da personalidade, sociedade da informação, proteção de dados, autonomia privada, identidade pessoal, consentimento.
    
    \textbf{Conceitos principais:}
    \begin{itemize}
        \item A evolução dos direitos da personalidade na sociedade informacional — do direito à privacidade ao direito à autodeterminação informativa.
        \item O consentimento como instrumento de tutela insuficiente diante das assimetrias de poder e informação no ambiente digital.
        \item A identidade pessoal como construção social ameaçada pela redução a perfis algorítmicos.
        \item A necessidade de repensar os direitos da personalidade à luz das novas tecnologias, sem perder sua fundamentação antropológica.
        \item O direito ao esquecimento e o direito à explicação como novas manifestações dos direitos da personalidade na era digital.
    \end{itemize}
    
    \textbf{Articulação:} A análise dos direitos da personalidade na sociedade informacional desenvolvida por Doneda fornece o quadro teórico para compreender os desafios regulatórios da IA examinados no capítulo quarto. Sua crítica ao consentimento como instrumento insuficiente de tutela dialoga com a necessidade, apontada pela \MH{}, de uma regulação ética que vá além da mera autonomia individual e considere o bem comum como critério orientador do desenvolvimento tecnológico.
\end{fichamento}

\begin{fichamento}
    \textbf{Ref.:} \cite{Maritain1947PCG}
    
    \textbf{Conceitos-chave:} pessoa, bem comum, dignidade, sociabilidade, individualidade, personalidade, transcendência.
    
    \textbf{Conceitos principais:}
    \begin{itemize}
        \item A distinção entre individualidade (aquilo que no homem é matéria, princípio de individuação e egoísmo) e personalidade (aquilo que no homem é espírito, abertura ao outro e ao absoluto).
        \item A pessoa como ``intelecto e vontade'', ser espiritual que transcende a ordem material e social.
        \item O bem comum como fim da sociedade que não se reduz à soma dos bens individuais, mas inclui a realização integral de cada pessoa.
        \item A sociabilidade constitutiva da pessoa — o homem é naturalmente social e se realiza na comunidade política.
        \item A subordinação do bem particular ao bem comum e a simultânea transcendência da pessoa em relação à sociedade política.
    \end{itemize}
    
    \textbf{Articulação:} A distinção entre individualidade e personalidade em Maritain fundamenta a tese da irredutibilidade da pessoa que a \MH{} sustenta diante do reducionismo algorítmico. Se a personalidade é abertura espiritual e transcendência, então qualquer sistema de IA que pretenda capturar a totalidade da pessoa através de dados e algoritmos comete um erro ontológico fundamental. O bem comum como fim da sociedade orienta a concepção de governança algorítmica proposta no capítulo quinto.
\end{fichamento}

\begin{fichamento}
    \textbf{Ref.:} \cite{Wojtyla1979PA}
    
    \textbf{Conceitos-chave:} pessoa, ação, consciência, transcendência, dignidade, autodeterminação, experiência vivida.
    
    \textbf{Conceitos principais:}
    \begin{itemize}
        \item A pessoa se manifesta através da ação — a ação é o meio privilegiado para conhecer a pessoa e sua estrutura ontológica.
        \item A transcendência da pessoa na ação — o ser humano não é apenas determinado por causas externas, mas se autodetermina através de escolhas livres e responsáveis.
        \item A consciência como elemento constitutivo da ação pessoal — a pessoa age com conhecimento e vontade, diferentemente de um mero processador de inputs.
        \item A autodeterminação como capacidade de agir sobre si mesmo, de se autogovernar, de ser causa de suas próprias ações.
        \item A dignidade da pessoa como sujeito moral — a pessoa é fins em si mesma, jamais meio para fins alheios.
    \end{itemize}
    
    \textbf{Articulação:} A fenomenologia da ação pessoal de Karol Wojtyła informa a compreensão da pessoa como ser livre e autorresponsável, que a \MH{} contrapõe à visão determinista dos algoritmos. A tese da transcendência na ação — o ser humano não se reduz a inputs e outputs — fundamenta a crítica radical da encíclica à objetificação algorítmica, que ignora a capacidade humana de autodeterminação e reduz a pessoa a um conjunto de comportamentos previsíveis e manipuláveis.
\end{fichamento}

\section{Artigos Científicos}

\begin{fichamento}
    \textbf{Ref.:} \cite{Floridi2018AIG}
    
    \textbf{Conceitos-chave:} ética da IA, princípios éticos, beneficência, não maleficência, autonomia, justiça, explicabilidade.
    
    \textbf{Conceitos principais:}
    \begin{itemize}
        \item Quatro princípios éticos para a IA: beneficência, não maleficência, autonomia e justiça.
        \item A explicabilidade como princípio transversal.
        \item A necessidade de traduzir princípios éticos em práticas regulatórias.
    \end{itemize}
    
    \textbf{Articulação:} Oferece um quadro comparativo para os princípios éticos da \MH{}, examinados no capítulo segundo e aplicados no capítulo quinto.
\end{fichamento}

\begin{fichamento}
    \textbf{Ref.:} \cite{Calo2017AIL}
    
    \textbf{Conceitos-chave:} IA, regulação, direito, mercado, segurança, responsabilidade.
    
    \textbf{Conceitos principais:}
    \begin{itemize}
        \item Desafios regulatórios específicos da IA.
        \item A inadequação dos modelos regulatórios tradicionais.
        \item A necessidade de regulação adaptativa.
    \end{itemize}
    
    \textbf{Articulação:} Fundamenta a análise das insuficiências do ordenamento jurídico brasileiro diante dos desafios da IA, no capítulo quarto.
\end{fichamento}

\begin{fichamento}
    \textbf{Ref.:} \cite{Mulholland2021IA}
    
    \textbf{Conceitos-chave:} inteligência artificial, direito, regulação, responsabilidade, direitos fundamentais, dignidade.
    
    \textbf{Conceitos principais:}
    \begin{itemize}
        \item Os desafios específicos da IA para o direito brasileiro: opacidade, imprevisibilidade e autonomia dos sistemas algorítmicos.
        \item A responsabilidade civil por danos causados por sistemas de IA — da responsabilidade subjetiva à responsabilidade objetiva e ao risco criado.
        \item A insuficiência dos institutos tradicionais do direito privado diante dos riscos sistêmicos da inteligência artificial.
        \item A necessidade de uma abordagem interdisciplinar que articule direito, ética e tecnologia para a regulação da IA.
        \item A proteção de dados pessoais como primeiro passo para a regulação da IA no Brasil.
    \end{itemize}
    
    \textbf{Articulação:} O artigo fundamenta a análise dos desafios jurídicos específicos da IA no direito brasileiro, desenvolvida no capítulo quarto. A proposta de uma abordagem interdisciplinar dialoga com a perspectiva integral da \MH{}, que reconhece a complexidade da relação entre técnica, ética e direito na proteção da personalidade humana.
\end{fichamento}

\begin{fichamento}
    \textbf{Ref.:} \cite{Wachter2017TSP}
    
    \textbf{Conceitos-chave:} direito à explicação, GDPR, decisões automatizadas, transparência, opacidade algorítmica, profiling.
    
    \textbf{Conceitos principais:}
    \begin{itemize}
        \item O direito à explicação sob o GDPR: alcance e limites do Art. 22 do Regulamento Europeu.
        \item A distinção entre explicação \textit{ex ante} (antes da decisão) e explicação \textit{ex post} (após a decisão) como graus diversos de proteção do titular.
        \item A insuficiência do direito à explicação nos casos de modelos algorítmicos complexos (\textit{black box}).
        \item A tensão entre o direito à explicação e os segredos comerciais protegidos pelas empresas de tecnologia.
        \item A proposta de um ``direito à informação significativa'' (\textit{right to meaningful information}) sobre o funcionamento dos sistemas algorítmicos.
    \end{itemize}
    
    \textbf{Articulação:} A defesa do direito à explicação como garantia fundamental nas decisões automatizadas fundamenta a análise da LGPD no capítulo quarto e dialoga com o princípio de transparência proposto pela \MH{}. As limitações identificadas pelas autoras informam a necessidade de um aperfeiçoamento regulatório que vá além do mero direito à explicação e assegure uma transparência substancial dos sistemas algorítmicos.
\end{fichamento}

\begin{fichamento}
    \textbf{Ref.:} \cite{Barocas2016BFF}
    
    \textbf{Conceitos-chave:} discriminação, big data, viés algorítmico, equidade, justiça, disparate impact.
    
    \textbf{Conceitos principais:}
    \begin{itemize}
        \item O impacto discriminatório dos sistemas de big data — mesmo sem intenção de discriminar, os algoritmos podem perpetuar e amplificar desigualdades existentes.
        \item A dificuldade de detectar discriminação algorítmica quando os critérios de decisão são opacos e os dados históricos já refletem preconceitos sociais.
        \item A distinção entre tratamento desigual (\textit{disparate treatment}) e impacto desproporcional (\textit{disparate impact}) como categorias analíticas para avaliar a justiça algorítmica.
        \item O ``ciclo de feedback'' entre dados discriminatórios e decisões algorítmicas — o viés se autoalimenta.
        \item A insuficiência das abordagens meramente técnicas para a equidade algorítmica — a justiça algorítmica exige também intervenção normativa.
    \end{itemize}
    
    \textbf{Articulação:} A análise dos vieses discriminatórios em sistemas automatizados desenvolvida por Barocas fundamenta a crítica da \MH{} à objetificação algorítmica da pessoa. O conceito de impacto desproporcional (\textit{disparate impact}) oferece um critério jurídico para avaliar a equidade dos sistemas de IA, complementar ao princípio de justiça defendido pela encíclica, informando as propostas regulatórias do capítulo quinto.
\end{fichamento}

\begin{fichamento}
    \textbf{Ref.:} \cite{Chesney2019DF}
    
    \textbf{Conceitos-chave:} desinformação, deep fakes, regulação, democracia, segurança, manipulação, autenticidade.
    
    \textbf{Conceitos principais:}
    \begin{itemize}
        \item As deep fakes como nova forma de desinformação que utiliza IA para criar conteúdos audiovisuais falsos indistinguíveis da realidade.
        \item Os riscos das deep fakes para a democracia: manipulação eleitoral, destruição de reputações e erosão da confiança pública.
        \item As deep fakes como ameaça à identidade pessoal — a possibilidade de atribuir falsamente atos e palavras a uma pessoa.
        \item Os desafios regulatórios específicos das deep fakes: detecção técnica insuficiente, liberdade de expressão e jurisdição transnacional.
        \item A necessidade de uma resposta combinada de tecnologia, direito e educação midiática para enfrentar o fenômeno.
    \end{itemize}
    
    \textbf{Articulação:} O fenômeno das deep fakes ilustra uma das ameaças concretas à personalidade que a \MH{} aborda indiretamente — a manipulação algorítmica que atinge a identidade, a honra e a imagem da pessoa. A análise dos riscos das deep fakes fundamenta a necessidade, defendida no capítulo quinto, de uma regulação específica que proteja a personalidade contra as formas mais invasivas de manipulação algorítmica.
\end{fichamento}

\begin{fichamento}
    \textbf{Ref.:} \cite{DoshiVelez2017MAR}
    
    \textbf{Conceitos-chave:} accountability, IA, responsabilidade, explicabilidade, prestação de contas, governança.
    
    \textbf{Conceitos principais:}
    \begin{itemize}
        \item A accountability (prestação de contas) como princípio fundamental para a governança dos sistemas de IA.
        \item A tensão entre a complexidade técnica dos modelos de IA e a exigibilidade de accountability pública.
        \item A necessidade de mecanismos de auditoria independente para sistemas algorítmicos que afetam direitos fundamentais.
        \item O direito à explicação como instrumento de accountability — a transparência algorítmica como condição para a responsabilização.
        \item A proposta de standards mínimos de accountability para diferentes níveis de risco dos sistemas de IA.
    \end{itemize}
    
    \textbf{Articulação:} O artigo fundamenta a necessidade de sistemas de IA accountable, compatível com a exigência de responsabilidade (\textit{accountability}) da \MH{}. A proposta de standards mínimos de accountability por nível de risco dialoga diretamente com a classificação de risco proposta pelo PL 2.338/2023 e analisada no capítulo quarto, oferecendo critérios para seu aperfeiçoamento à luz dos princípios éticos da encíclica.
\end{fichamento}

\begin{fichamento}
    \textbf{Ref.:} \cite{DoshiVelez2017TMI}
    
    \textbf{Conceitos-chave:} IA interpretável, explicabilidade, rigor científico, complexidade, transparência, confiança.
    
    \textbf{Conceitos principais:}
    \begin{itemize}
        \item A distinção entre modelos intrinsecamente interpretáveis (como árvores de decisão) e modelos pós-hoc explicáveis (explicações externas a modelos complexos).
        \item A tensão fundamental entre acurácia e interpretabilidade — modelos mais precisos tendem a ser menos interpretáveis.
        \item A necessidade de um programa de pesquisa rigoroso para a IA explicável (\textit{explainable AI — XAI}).
        \item A identificação de diferentes públicos para a explicação (usuários finais, reguladores, desenvolvedores) e suas diferentes necessidades informacionais.
        \item Os limites da explicabilidade — nem toda decisão algorítmica pode ser explicada de forma completa e compreensível.
    \end{itemize}
    
    \textbf{Articulação:} A exigência de IA interpretável fundamenta o direito à explicação analisado na pesquisa. A distinção entre explicabilidade intrínseca e pós-hoc informa a análise crítica dos limites do direito à explicação na LGPD e no PL 2.338/2023, desenvolvida no capítulo quarto. A complexidade identificada pelos autores reforça a necessidade de uma abordagem normativa que a \MH{} oferece: a transparência algorítmica como exigência ética, não apenas técnica.
\end{fichamento}

\begin{fichamento}
    \textbf{Ref.:} \cite{Edwards2017PC}
    
    \textbf{Conceitos-chave:} direito à explicação, GDPR, profiling, opacidade, autonomia, decisões automatizadas.
    
    \textbf{Conceitos principais:}
    \begin{itemize}
        \item A análise crítica do direito à explicação no GDPR: o alcance limitado do Art. 22 e sua baixa efetividade prática.
        \item A distinções entre diferentes tipos de decisões automatizadas (perfilamento, categorização, decisões individuais) e seus diferentes regimes de proteção.
        \item A opacidade dos sistemas de profiling como ameaça à autonomia pessoal — a pessoa é categorizada e julgada sem ter conhecimento ou controle sobre o processo.
        \item A proposta de um direito à intervenção humana significativa, não apenas formal, nas decisões automatizadas.
        \item A crítica à abordagem europeia de regulação da IA como insuficiente para proteger os direitos fundamentais dos cidadãos.
    \end{itemize}
    
    \textbf{Articulação:} A defesa do direito à explicação como garantia fundamental, aliada à crítica de suas limitações práticas no GDPR, fundamenta a análise da LGPD e dos projetos de lei brasileiros no capítulo quarto. A proposta de intervenção humana significativa dialoga com a exigência da \MH{} de que a técnica esteja sempre subordinada ao juízo ético da pessoa, jamais o substituindo de forma irreversível.
\end{fichamento}

\begin{fichamento}
    \textbf{Ref.:} \cite{Goodman2017EU}
    
    \textbf{Conceitos-chave:} regulação, tomada de decisão algorítmica, GDPR, direitos, transparência, não discriminação.
    
    \textbf{Conceitos principais:}
    \begin{itemize}
        \item A análise das disposições do GDPR relevantes para a tomada de decisão algorítmica, especialmente os Arts. 22, 13, 14 e 15.
        \item O direito à explicação como inovação do GDPR e seus fundamentos na tradição jurídica europeia de proteção de dados.
        \item A proibição de decisões exclusivamente automatizadas como regra geral, com exceções limitadas.
        \item As obrigações de transparência do controlador: informação significativa sobre a lógica do algoritmo.
        \item A avaliação crítica da eficácia do GDPR diante da complexidade técnica dos modelos de aprendizado de máquina.
    \end{itemize}
    
    \textbf{Articulação:} O artigo oferece a base comparativa europeia para a análise do regime jurídico brasileiro de regulação da IA no capítulo quarto. A estrutura de direitos e garantias do GDPR — especialmente o direito à explicação e a proibição de decisões exclusivamente automatizadas — constitui o parâmetro mínimo de proteção contra o qual a LGPD e o PL 2.338/2023 são avaliados, em diálogo com os princípios éticos propostos pela \MH{}.
\end{fichamento}

\begin{fichamento}
    \textbf{Ref.:} \cite{Binns2018ITG}
    
    \textbf{Conceitos-chave:} equidade, aprendizado de máquina, justiça algorítmica, imparcialidade, vieses, fairness.
    
    \textbf{Conceitos principais:}
    \begin{itemize}
        \item As diferentes definições formais de equidade (\textit{fairness}) no aprendizado de máquina: igualdade de oportunidades, paridade demográfica, imparcialidade individual.
        \item A incompatibilidade entre diferentes noções de equidade — não é possível satisfazer simultaneamente todos os critérios formais de justiça algorítmica.
        \item A necessidade de escolhas normativas explícitas sobre qual concepção de equidade orienta o desenvolvimento de cada sistema de IA.
        \item A tensão entre equidade algorítmica e acurácia preditiva — sistemas mais justos podem ser menos precisos e vice-versa.
        \item A importância do contexto social e institucional na determinação do que constitui uma decisão algorítmica justa.
    \end{itemize}
    
    \textbf{Articulação:} O artigo fornece critérios técnicos e conceituais para avaliar a equidade dos sistemas de IA, compatível com o princípio de justiça da \MH{}. A demonstração de que não existe uma única definição neutra de equidade algorítmica reforça a exigência da encíclica de que o desenvolvimento tecnológico seja orientado por uma visão integral da pessoa humana e do bem comum, e não por critérios meramente formais ou procedimentais.
\end{fichamento}

\begin{fichamento}
    \textbf{Ref.:} \cite{Selbst2018INC}
    
    \textbf{Conceitos-chave:} equidade, intuição, discricionariedade, justiça algorítmica, normas sociais, contexto.
    
    \textbf{Conceitos principais:}
    \begin{itemize}
        \item A tensão entre as noções intuitivas de equidade (\textit{intuitive notions of fairness}) e as definições formais de equidade algorítmica.
        \item A discricionariedade humana como elemento que sistemas automatizados eliminam, mas que pode ser necessária para decisões justas e contextualizadas.
        \item A diferença entre a flexibilidade interpretativa humana (que permite adaptar regras a casos específicos) e a rigidez dos sistemas algorítmicos.
        \item A perda de nuances contextuais na tradução de decisões sociais complexas para modelos matemáticos.
        \item A necessidade de preservar espaço para o juízo humano e a discricionariedade informada em decisões com impacto significativo sobre direitos.
    \end{itemize}
    
    \textbf{Articulação:} As limitações da equidade algorítmica identificadas por Selbst — especialmente a perda de discricionariedade contextual — informam a necessidade de uma abordagem normativa mais ampla como a proposta pela \MH{}. A crítica à rigidez dos sistemas automatizados reforça a defesa da encíclica pela irredutibilidade da pessoa a critérios algorítmicos e pela necessidade de preservar o juízo ético humano como instância última de decisão.
\end{fichamento}

\begin{fichamento}
    \textbf{Ref.:} \cite{Selbst2018TIR}
    
    \textbf{Conceitos-chave:} regulação, insuficiência, adequação, direitos, transparência, contexto.
    
    \textbf{Conceitos principais:}
    \begin{itemize}
        \item A inadequação dos marcos regulatórios existentes para lidar com os desafios específicos dos sistemas de IA.
        \item Os ``desalinhamentos'' (\textit{mismatches}) entre as características técnicas da IA e os pressupostos dos modelos regulatórios tradicionais.
        \item A identificação de cinco categorias de desalinhamento: tecnológico, sociotécnico, de agência, de conhecimento e de contexto.
        \item A necessidade de reformular as abordagens regulatórias para considerar a natureza sociotécnica dos sistemas algorítmicos.
        \item A proposta de uma regulação contextual e adaptativa, que reconheça as especificidades de cada domínio de aplicação da IA.
    \end{itemize}
    
    \textbf{Articulação:} A tese da inadequação dos marcos regulatórios existentes — em particular os cinco desalinhamentos identificados por Selbst — fundamenta a análise crítica do capítulo quarto sobre as insuficiências do ordenamento brasileiro diante dos desafios da IA. A proposta de regulação contextual dialoga com a abordagem da \MH{}, que também reconhece a necessidade de princípios éticos que orientem a aplicação da técnica em contextos específicos, sem cair em generalizações abstratas ou prescrições inflexíveis.
\end{fichamento}

\begin{fichamento}
    \textbf{Ref.:} \cite{Bostrom2005TH}
    
    \textbf{Conceitos-chave:} transumanismo, superinteligência, ética, futuro da humanidade, aperfeiçoamento humano, risco existencial.
    
    \textbf{Conceitos principais:}
    \begin{itemize}
        \item O transumanismo como movimento intelectual e cultural que defende o aperfeiçoamento da condição humana através da tecnologia.
        \item A superinteligência como objetivo último do transumanismo — uma inteligência artificial que supera o melhor cérebro humano em todos os domínios.
        \item Os riscos existenciais associados ao desenvolvimento descontrolado da IA — a possibilidade de extinção da humanidade.
        \item A crítica às objeções bioéticas e humanistas ao transumanismo, que Bostrom considera moralmente paralisantes.
        \item A defesa do aperfeiçoamento técnico do ser humano como extensão do projeto humanista de autossuperação.
    \end{itemize}
    
    \textbf{Articulação:} A análise do transumanismo por Bostrom — que defende o aperfeiçoamento humano através da tecnologia — informa a crítica da \MH{} ao paradigma tecnocrático. A defesa bostromiana da superação dos limites humanos pela técnica representa a visão antagônica à antropologia integral da encíclica, que afirma a irredutibilidade da pessoa à sua dimensão biológica ou informacional. O contraste entre as duas visões, desenvolvido no capítulo terceiro, permite demonstrar a originalidade e a atualidade da proposta ética da \MH{}.
\end{fichamento}

\section{Legislação e Documentos}

\begin{fichamento}
    \textbf{Ref.:} \cite{Brasil1988CF}
    
    \textbf{Conceitos-chave:} dignidade da pessoa humana, direitos fundamentais, personalidade, privacidade, honra, imagem, inviolabilidade.
    
    \textbf{Dispositivos relevantes:}
    \begin{itemize}
        \item Art. 1º, III: Dignidade da pessoa humana como fundamento da República.
        \item Art. 5º, X: Inviolabilidade da vida privada, honra e imagem.
        \item Art. 5º, LXXII: Habeas data como garantia fundamental.
    \end{itemize}
    
    \textbf{Articulação:} A Constituição Federal constitui o fundamento normativo máximo da proteção da personalidade, examinada no capítulo quarto em seus limites e possibilidades diante dos desafios da IA.
\end{fichamento}

\begin{fichamento}
    \textbf{Ref.:} \cite{Brasil2018LGPD}
    
    \textbf{Conceitos-chave:} proteção de dados, consentimento, direito à explicação, decisões automatizadas, ANPD, sanções.
    
    \textbf{Dispositivos relevantes:}
    \begin{itemize}
        \item Art. 6º: Princípios da proteção de dados (finalidade, adequação, necessidade, transparência, prevenção, não discriminação).
        \item Art. 20: Direito à revisão de decisões automatizadas.
        \item Art. 38: Obrigação de comunicação de incidentes de segurança.
    \end{itemize}
    
    \textbf{Articulação:} A LGPD é o principal instrumento jurídico brasileiro de proteção de dados pessoais. Sua análise crítica no capítulo quarto e a articulação com os princípios da \MH{} no capítulo quinto constituem o núcleo normativo da pesquisa.
\end{fichamento}

\begin{fichamento}
    \textbf{Ref.:} \cite{Brasil2023PL2338}
    
    \textbf{Conceitos-chave:} IA, regulação, classificação de risco, transparência, governança, accountability, fiscalização.
    
    \textbf{Dispositivos relevantes:}
    \begin{itemize}
        \item Classificação dos sistemas de IA por nível de risco.
        \item Obrigações de transparência e explicabilidade.
        \item Criação de autoridade competente para fiscalização.
    \end{itemize}
    
    \textbf{Articulação:} O PL 2.338/2023 representa a tentativa mais recente de regulação específica da IA no Brasil. Sua análise no capítulo quarto e a articulação com os princípios éticos da \MH{} no capítulo quinto apontam caminhos para seu aperfeiçoamento.
\end{fichamento}

\begin{fichamento}
    \textbf{Ref.:} \cite{Brasil2002CC}
    
    \textbf{Conceitos-chave:} pessoa natural, capacidade, direitos da personalidade, nome, imagem, honra, privacidade, disposição do próprio corpo.
    
    \textbf{Dispositivos relevantes:}
    \begin{itemize}
        \item Arts. 1º–10: Da pessoa natural — capacidade jurídica, nascimento, morte, ausência, domicílio.
        \item Arts. 11–21: Dos direitos da personalidade — intransmissibilidade, irrenunciabilidade, tutela do nome, imagem, honra e privacidade.
        \item Art. 12: Ação de cessação e indenização por violação a direitos da personalidade, inclusive após a morte.
        \item Art. 20: Proteção da imagem e da palavra — proibição de divulgação de escritos, palavras ou imagem sem autorização.
        \item Arts. 186–187: Responsabilidade civil por ato ilícito e abuso de direito como fundamento para reparação de danos à personalidade causados por sistemas de IA.
    \end{itemize}
    
    \textbf{Articulação:} O Código Civil fundamenta a proteção dos direitos da personalidade no direito civil brasileiro, examinada no capítulo quarto. A abertura do sistema de direitos da personalidade — que a doutrina reconhece como rol exemplificativo — permite a tutela de novas manifestações da personalidade ameaçadas pela IA, como a autodeterminação informativa e o direito à explicação algorítmica, em harmonia com a visão integral da pessoa proposta pela \MH{}.
\end{fichamento}

\begin{fichamento}
    \textbf{Ref.:} \cite{UE2016GDPR}
    
    \textbf{Conceitos-chave:} proteção de dados, privacidade, GDPR, decisões automatizadas, direito à explicação, transferência internacional, consentimento.
    
    \textbf{Dispositivos relevantes:}
    \begin{itemize}
        \item Arts. 13–14: Obrigações de transparência do controlador — informações sobre a lógica do tratamento e as consequências previsíveis para o titular.
        \item Art. 22: Direito do titular a não ficar sujeito a decisões exclusivamente automatizadas, incluindo perfilamento, com exceções limitadas e garantias.
        \item Arts. 35–36: Obrigação de realizar avaliação de impacto sobre a proteção de dados (\textit{Data Protection Impact Assessment}) para tratamentos de alto risco.
        \item Art. 25: Proteção de dados desde a concepção e por padrão (\textit{privacy by design and by default}).
        \item Art. 46: Transferência internacional de dados pessoais — garantias adequadas e decisões de adequação.
    \end{itemize}
    
    \textbf{Articulação:} O GDPR é o marco regulatório de referência internacional para proteção de dados, que influencia a LGPD e o PL 2.338/2023 no Brasil. Sua análise comparativa no capítulo quarto permite identificar tanto os avanços da legislação brasileira quanto suas lacunas face ao direito europeu, enquanto os princípios de transparência, privacidade desde a concepção e avaliação de impacto dialogam com a exigência da \MH{} de subordinar o desenvolvimento técnico ao controle ético e à proteção integral da pessoa humana.
\end{fichamento}

\begin{fichamento}
    \textbf{Ref.:} \cite{STF2020ADI6640}
    
    \textbf{Conceitos-chave:} proteção de dados, direito fundamental, ADI 6640, STF, LGPD, dignidade, privacidade.
    
    \textbf{Dispositivos relevantes:}
    \begin{itemize}
        \item Voto do relator, Ministro Luiz Fux: A proteção de dados pessoais como direito fundamental autônomo, extraído do princípio da dignidade da pessoa humana e do direito à privacidade.
        \item Tese firmada: A LGPD é constitucional e a proteção de dados é direito fundamental com eficácia imediata, nos termos do Art. 5º, §2º da CF.
        \item Fundamento: A autodeterminação informativa como projeção do livre desenvolvimento da personalidade e da autonomia privada.
        \item Reconhecimento da competência concorrente da União para legislar sobre proteção de dados, superando controvérsias federativas.
        \item Precedente de modulação: O STF estabeleceu o prazo de 24 meses para adequação à LGPD, reconhecendo a complexidade da transição.
    \end{itemize}
    
    \textbf{Articulação:} O STF, na ADI 6640, reconheceu a proteção de dados como direito fundamental autônomo, fundamentando a necessidade de regulação da IA no Brasil. O voto do relator, que conecta a proteção de dados à dignidade da pessoa humana e ao livre desenvolvimento da personalidade, dialoga diretamente com a fundamentação antropológica da \MH{}, que também vincula a proteção da pessoa na era digital à sua dignidade intrínseca e irredutível.
\end{fichamento}

\begin{fichamento}
    \textbf{Ref.:} \cite{STF2022ADPF695}
    
    \textbf{Conceitos-chave:} compartilhamento de dados, administração pública, privacidade, proteção de dados, limites ao poder estatal, LGPD.
    
    \textbf{Dispositivos relevantes:}
    \begin{itemize}
        \item Decisão do STF que estabeleceu limites ao compartilhamento de dados entre órgãos da administração pública federal.
        \item Exigência de base legal específica e de transparência para o compartilhamento de dados pessoais entre entidades públicas.
        \item Reafirmação da aplicabilidade da LGPD ao setor público, inclusive no que tange às sanções por violação.
        \item Necessidade de avaliação de impacto à proteção de dados antes de implementar sistemas de compartilhamento em larga escala.
        \item Vinculação do tratamento de dados pela administração pública aos princípios constitucionais da impessoalidade, moralidade e eficiência.
    \end{itemize}
    
    \textbf{Articulação:} A decisão do STF na ADPF 695 reafirmou limites à utilização de dados pela administração pública, tema relevante para a regulação da IA no setor público, examinada no capítulo quarto. O reconhecimento de que o poder público está vinculado à LGPD e aos princípios constitucionais quando trata dados pessoais dialoga com a exigência da \MH{} de que o desenvolvimento tecnológico seja orientado pelo bem comum e pelo respeito à dignidade de cada pessoa, especialmente aquelas em situação de vulnerabilidade.
\end{fichamento}

\begin{fichamento}
    \textbf{Ref.:} \cite{HighLevelExpertGroup2019EA}
    
    \textbf{Conceitos-chave:} IA confiável, licitude, eticidade, robustez técnica, transparência, accountability, bem-estar social.
    
    \textbf{Dispositivos relevantes:}
    \begin{itemize}
        \item Cap. I: Requisitos para uma IA confiável — ação humana e supervisão, robustez técnica e segurança, privacidade e governança de dados, transparência, diversidade e não discriminação, bem-estar social e ambiental, accountability.
        \item Cap. II: Avaliação de confiabilidade — lista de verificação (\textit{assessment list}) para desenvolvedores e implementadores de sistemas de IA.
        \item Cap. III: Governança da IA — necessidade de mecanismos de governança multinível que articulem autorregulação, regulação estatal e participação social.
        \item Recomendações para políticas públicas: investimento em pesquisa, educação digital, fortalecimento de autoridades reguladoras e cooperação internacional.
        \item Ênfase na centralidade da pessoa (\textit{human-centric AI}) como princípio orientador de todo desenvolvimento tecnológico.
    \end{itemize}
    
    \textbf{Articulação:} As diretrizes éticas do Grupo de Peritos de Alto Nível da UE constituem o principal documento internacional de ética da IA, analisado comparativamente com a \MH{} no capítulo segundo. A centralidade da pessoa (\textit{human-centric AI}) como princípio orientador converge com a visão antropológica da encíclica, enquanto os sete requisitos para IA confiável oferecem um quadro operacional que pode ser enriquecido pelos princípios éticos mais amplos propostos pela \MH{}, especialmente o princípio da irredutibilidade da pessoa e o destino universal dos bens aplicado aos dados.
\end{fichamento}

